\section{2016年高考题}
\subsection{2016年全国I卷（理）}\label{gk:2016年全国I卷（理）}


\clearpage


\subsection{2016年全国I卷（文）}\label{gk:2016年全国I卷（文）}


\clearpage


\subsection{2016年全国II卷（理）}\label{gk:2016年全国II卷（理）}


\clearpage


\subsection{2016年全国II卷（文）}\label{gk:2016年全国II卷（文）}


\clearpage


\subsection{2016年全国III卷（理）}\label{gk:2016年全国III卷（理）}


\clearpage


\subsection{2016年全国III卷（文）}\label{gk:2016年全国III卷（文）}


\clearpage


\subsection{2016年北京卷（理）}\label{gk:2016年北京卷（理）}

\begin{description}
    \item[一、]
    \begin{enumerate}
        \item 1
        \item 2
        \item 3
        \item 4
        \item 5
        \item 6
        \item 7
        \item 8
    \end{enumerate}
    \item[二、]
    \begin{enumerate}
      \addtocounter{enumi}{+8}   
        \item 1
        \item 2
        \item 3
        \item 4
        \item 1
        \item 6
    \end{enumerate}
    \item[三、]   
    \begin{enumerate}
      \addtocounter{enumi}{+14}   
        \item 1
        \item 2
        \item 3
        \item 设函数 $f(x) = x e^{a - x} + b x$，曲线 $y = f(x)$ 在点 $(2, f(2))$ 处的切线方程为 $y = (e - 1)x + 4$。
\begin{enumerate}[label=(\arabic*)]
    \item 求 $a$, $b$ 的值；
    \item 求 $f(x)$ 的单调区间。
\end{enumerate}

        \item 已知椭圆 $C: \frac{x^{2}}{a^{2}} + \frac{y^{2}}{b^{2}} = 1$ ($a > b > 0$) 的离心率为 $\frac{\sqrt{3}}{2}$，$A(a,0)$，$B(0,b)$，$O(0,0)$，$\triangle OAB$ 的面积为 $1$。

\begin{enumerate}[label=(\arabic*)]
    \item 求椭圆 $C$ 的方程；
    \item 设 $P$ 是椭圆 $C$ 上一点，直线 $PA$ 与 $y$ 轴交于点 $M$，直线 $PB$ 与 $x$ 轴交于点 $N$。求证：$|AN| \cdot |BM|$ 为定值。
\end{enumerate}
        \item 设数列 $A: a_{1}, a_{2}, \cdots, a_{N}$ ($N \geqslant 2$)。如果对小于 $n$ ($2 \leqslant n \leqslant N$) 的每个正整数 $k$ 都有 $a_{k} < a_{n}$，则称 $n$ 是数列 $A$ 的一个"G 时刻"。记 $G(A)$ 是数列 $A$ 的所有"G 时刻"组成的集合。

\begin{enumerate}[label=(\arabic*)]
    \item 对数列 $A: -2, 2, -1, 1, 3$，写出 $G(A)$ 的所有元素；
    \item 证明：若数列 $A$ 中存在 $a_{n}$ 使得 $a_{n} > a_{1}$，则 $G(A) \neq \emptyset$；
    \item 证明：若数列 $A$ 满足 $a_{n} - a_{n-1} \leqslant 1$ ($n = 2, 3, \cdots, N$)，则 $G(A)$ 的元素个数不小于 $a_{N} - a_{1}$。
\end{enumerate}
    \end{enumerate}
\end{description}
\clearpage


\subsection{2016年北京卷（文）}\label{gk:2016年北京卷（文）}

\begin{description}
    \item[一、]
    \begin{enumerate}
        \item 1
        \item 2
        \item 3
        \item 4
        \item 5
        \item 6
        \item 7
        \item 8
    \end{enumerate}
    \item[二、]
    \begin{enumerate}
      \addtocounter{enumi}{+8}   
        \item 1
        \item 2
        \item 3
        \item 4
        \item 1
        \item 6
    \end{enumerate}
    \item[三、]   
    \begin{enumerate}
      \addtocounter{enumi}{+14}   
        \item 1
        \item 2
        \item 3
        \item 
        \item 已知椭圆 $C: \frac{x^{2}}{a^{2}} + \frac{y^{2}}{b^{2}} = 1$ 过点 $A(2,0)$，$B(0,1)$ 两点。

\begin{enumerate}[label=(\arabic*)]
    \item 求椭圆 $C$ 的方程及离心率；
    \item 设 $P$ 为第三象限内一点且在椭圆 $C$ 上，直线 $PA$ 与 $y$ 轴交于点 $M$，直线 $PB$ 与 $x$ 轴交于点 $N$，求证：四边形 $ABNM$ 的面积为定值。
\end{enumerate}
        \item  设函数 $f(x) = x^{3} + a x^{2} + b x + c$。

\begin{enumerate}[label=(\arabic*)]
    \item 求曲线 $y = f(x)$ 在点 $(0, f(0))$ 处的切线方程；
    \item 设 $a = b = 4$，若函数 $f(x)$ 有三个不同零点，求 $c$ 的取值范围；
    \item 求证：$a^{2} - 3b > 0$ 是 $f(x)$ 有三个不同零点的必要而不充分条件。
\end{enumerate}
    \end{enumerate}
\end{description}
\clearpage


\subsection{2016年江苏卷}\label{gk:2016年江苏卷}


\clearpage


\subsection{2016年上海卷（理）}\label{gk:2016年上海卷（理）}


\clearpage


\subsection{2016年上海卷（文）}\label{gk:2016年上海卷（文）}


\clearpage


\subsection{2016年四川卷（理）}\label{gk:2016年四川卷（理）}


\clearpage


\subsection{2016年四川卷（文）}\label{gk:2016年四川卷（文）}


\clearpage


\subsection{2016年山东卷（理）}\label{gk:2016年山东卷（理）}


\clearpage


\subsection{2016年山东卷（文）}\label{gk:2016年山东卷（文）}


\clearpage


\subsection{2016年天津卷（理）}\label{gk:2016年天津卷（理）}


\clearpage


\subsection{2016年天津卷（文）}\label{gk:2016年天津卷（文）}


\clearpage


\subsection{2016年浙江卷（理）}\label{gk:2016年浙江卷（理）}


\clearpage


\subsection{2016年浙江卷（文）}\label{gk:2016年浙江卷（文）}


\clearpage

